\section{Introduction}
\label{sec:intro}


Mathematical Programs with Equilibrium Constraints
(MPECs), equivalently referred to as Mathematical Programs with Complementarity Constraints (MPCCs),\footnote{While the literature often uses MPEC and MPCC interchangeably, we adopt the term MPEC throughout this paper for consistency.} are optimization problems~\citep{LuoPangRalph1996,OutrataKocvaraZowe1998,ScheelScholtes2000} that arise in a wide range of applications, including bilevel optimization~\cite{Dempe2002,ColsonMarcotteSavard2007}, noncooperative game theory~\cite{LeyfferMunson2010,HuRalph2007}, contact mechanics~\cite{OutrataKocvaraZowe1998, KlarbringBjorkman1988}, traffic flow equilibrium~\cite{FerrisPang1997,MengYangBell2001}, nonsmooth optimal control~\cite{NurkanovicPozharskiyDiehl2024}, and chemical engineering process optimization~\cite{BaumruckerBiegler2008}.

Luo et al.~\cite{LuoPangRalph1996} established the formal
study of MPECs and formulated their foundational theory.
% Scheel and Scholtes~\cite{ScheelScholtes2000} subsequently formalized
% the MPEC stationarity hierarchy---introducing the concepts of W-, C-, M-,
% and S-stationarity---to characterize optimality in the absence of standard
% constraint qualifications (such as LICQ and MFCQ), which necessarily fail
% at every feasible point of an MPEC, rendering classical KKT theory
% inapplicable. 
In general, a standard MPEC can be formulated as follows~\cite{Ye2005}:
\begin{equation}
\begin{aligned}
  &\min_{x \in \mathbb{R}^n} \quad && f(x) \\
  &\text{subject to} \quad && g_i(x) \leq 0, \quad i = 1, \ldots, m, \\
  & && h_j(x) = 0, \quad j = 1, \ldots, p, \\
  & && G_\ell(x) \geq 0, \quad H_\ell(x) \geq 0, \quad \ell = 1, \ldots, q, \\
  & && G_\ell(x)^{\top} H_\ell(x) = 0, \quad \ell = 1, \ldots, q.
\end{aligned}
\label{eq:general_mpec}
\end{equation}
Here, $f(x)$ is the objective function, $g_i(x)$ and $h_j(x)$ represent
the standard inequality and equality constraints respectively.

The fundamental difficulty arising from complementarity constraints is
that standard constraint qualifications such as Linear Independence Constraint
Qualification (LICQ) and Mangasarian--Fromovitz Constraint Qualification
(MFCQ) are violated at all feasible points~\citep{ScheelScholtes2000}.
Consequently, classical KKT condition is failed, and
standard Nonlinear Programming (NLP) solvers applied directly
to~\eqref{eq:general_mpec} will not find valid solutions. 
To overcome this computational barrier several specialized algorithms have been developed over the past two decades. Notable developments include regularization
techniques~\cite{Scholtes2001,LinFukushima2005,Kadrani2009,SteffensenUlbrich2010}, sequential quadratic programming (SQP)
approaches~\cite{FletcherLeyffer2004,FletcherLeyfferRalphScholtes2006}, penalty
methods~\cite{HuRalph2004,Anitescu2005}, and interior-point
strategies~\cite{RaghunnathanBiegler2005,LeyfferLopezCalvaNocedal2006}.

% Several stationary points have been developed, among
% these, \emph{Bouligand (B-) stationarity} is particularly meaningful
% from a practical perspective: a B-stationary point admits no feasible
% descent direction in the tangent cone sense, making it a genuine
% first-order necessary condition for local optimality. Stronger (S-) stationarity is
% equivalent to B-stationarity under MPEC-LICQ and can be verified by 
% a fast algebraic test on multiplier signs, making it practical to 
% certify strong optimality at each iteration. The
% weaker \emph{Clarke (C-) stationarity} places no restriction on the
% signs of biactive complementarity multipliers, and C-stationary points
% can include saddle points or even local maxima that allow feasible
% descent directions. Standard solving methods often find weaker C-stationary points that are
% not genuinely optimal~\cite{HoheiselKanzowSchwartz2013}, making proper
% B-stationarity a vital algorithmic objective.

The most widely used approach for solving MPECs is \emph{Scholtes
regularization}~\cite{Scholtes2001}, which replaces the hard
complementarity product $G_\ell(x) \cdot H_\ell(x) = 0$ with the
relaxed condition $G_\ell(x) \cdot H_\ell(x) \leq t_k$, and drives
$t_k \to 0$ by solving a sequence of smooth NLP subproblems. While
numerically robust and widely applicable, the standard formulation
reduces $t_k$ via a fixed geometric rule that is independent of solver
progress, and is theoretically guaranteed to converge only to
C-stationary points~\cite{Scholtes2001,HoheiselKanzowSchwartz2013},
providing no built-in mechanism for targeting B-stationary. A detailed review of Scholtes regularization is given
in Section~\ref{sec:scholtes}; see also Hoheisel
et~al.~\cite{HoheiselKanzowSchwartz2013} for a comprehensive
theoretical and numerical comparison among different relaxation methods.

Existing methods face practical limitations: SQP approaches~\cite{FletcherLeyffer2004}
often depend on commercial solvers like SNOPT~\cite{GillMurraySaunders2005SNOPT},
and many practically competitive implementations still rely on
commercial linear-algebra or quadratic-programming components. The
recent MPECopt algorithm of Nurkanovi\'{c} and
Leyffer~\cite{NurkanovicLeyffer2025} is an open-source method with
global convergence to B-stationary points and strong reported
performance on the MacMPEC benchmark~\cite{Leyffer2000}; however, its
current implementation depends on the commercial MATLAB environment and
its strongest reported results use the commercial solver Gurobi.

Significant gaps remain in the practical MPEC software ecosystem. Many
algorithms rely on commercial software (HSL, SNOPT, Gurobi), standard
Scholtes regularization lacks built-in mechanisms for achieving
stronger stationarity than C-stationarity. MPECSS is designed to close these implementation and reproducibility gaps within a single open-source code base.

MPECSS uses the Scholtes framework while adding an implementation-aware
\emph{adaptive regularization update} driven by the
complementarity residual and an integrated \emph{multiplier sign test}
derived from the mathematical characterizations of Scheel and
Scholtes~\cite{ScheelScholtes2000}. In the implementation, this sign
test is used as an S-stationarity test on the $G$-multipliers and
$H$-multipliers at numerically biactive indices; final B-stationarity
certification is then performed in Phase~III by a combination of
MPEC-LICQ checks and, when needed, LPEC-based verification. The method
is implemented entirely in Python using the open-source CasADi
framework~\cite{AnderssonGillisHornRawlingsDiehl2019} and
IPOPT~\cite{WachterBiegler2006}, with no commercial dependencies.

The main contributions of this paper are summarized as follows:
\emph{MPECSS algorithm:} A fully open-source, multi-phase
  algorithm for solving MPECs that combines Scholtes regularization with
  adaptive parameter control, a multiplier-based stationarity
  test, multistart and restoration mechanisms, and a Phase~III
  certification pipeline combining LICQ checks, BNLP polishing, and
  LPEC refinement.
\emph{Adaptive regularization control:} A monitoring-driven update of
  $t_k$ based on the implemented homotopy residual, stable biactive-set
  tracking, and stagnation counters, together with practical safeguards
  such as adaptive jumps, jump caps, and a floor on $t_k$.
\emph{Multiplier sign test:} A fast algebraic procedure that checks the
  converted $G$-multipliers and $H$-multipliers on numerically biactive indices to
  detect S-stationarity candidates and steer the subsequent
  certification logic.
\emph{Open-source benchmark infrastructure:} A reproducible benchmark
  runner for MacMPEC, MPECLib, and NOSBENCH with machine-generated,
  wide-schema CSV logging (300+ columns),
  environment capture, process-level isolation for per-problem
  timeout enforcement, resume support (via \texttt{--resume}),
  retry-failed logic, and graceful recovery from timeout, interrupt,
  and out-of-memory events. The solver stack auto-selects between
  SQP$+$qpOASES (for problems with $n \leq 400$ when available) and
  IPOPT$+$MUMPS (otherwise).
\emph{Rigorous convergence guarantees:} We provide formal proofs that
  MPECSS inherits the baseline C-stationarity guarantees of Scholtes
  regularization, while stronger B-stationarity certificates are
  produced in Phase~III when LICQ equivalence or LPEC verification
  succeeds.


